\startsection[title={Fermentação de hidrogênio}]
	Os combustíveis fósseis possibilitaram avanços significativos para o desenvolvimento da sociedade moderna ao longo dos séculos, contribuindo para a consolidação da industrialização, da expansão dos meios de transporte e da intensificação da globalização. Embora a implementação inicial de métodos de geração energética em larga escala tenha sido fundamental para o desenvolvimento socioeconômico da humanidade, grande parte desses modelos apresentam impactos expressivos sobre a manutenção da estabilidade ambiental e da saúde do planeta a longo prazo. Atualmente, considerando o avanço tecnológico e a viabilidade de implementação de alternativas ao modelo energético predominantemente utilizado \cite[authoryear][ref::nanda_0], torna-se relevante destacar os avanços científicos relacionados a métodos alternativos de geração de energia renovaveis. Dentre elas, as rotas de conversão de resíduos orgânicos por meio de tratamentos termoquímicos e bioquímicos destacam-se como algumas das opções mais promissoras em substituição aos combustíveis fósseis \cite[authoryear][ref::nanda_0].

	O hidrogênio possui elevado potencial para utilização em larga escala como combustível. Nesse contexto, pode ser empregado diretamente como um biocombustível gasoso ou convertido em hidrocarbonetos e eletricidade por meio de células a combustível \cite[authoryear][ref::das_0]. Para fins de combustão, o hidrogênio apresenta valor calorífico inferior de aproximadamente \math{120 \unit{MJ} \unit{kg}^{-1}} \cite[authoryear][ref::sarangi_0], equivalente ao fornecido por cerca de \math{2.75 \unit{kg}} de gasolina \cite[authoryear][ref::nanda_0]. Tais características, associadas ao fato de sua combustão gerar água como único produto, tornam o hidrogênio uma alternativa energeticamente vantajosa, promissora e sustentável.

	As rotas para produção de hidrogênio são diversas, podendo ter como matéria-prima hidrocarbonetos, gás natural, água ou biomassa. Entre os principais processos, destacam-se a reforma de gás natural, decomposição térmica do gás natural, gaseificação de carvão, oxidação parcial de hidrocarbonetos, pirólise, eletrólise, fotólise, termólise e as rotas biológicas \cite[authoryear][ref::das_0]. Essas rotas, predominantemente termoquímicas e eletroquímicas, em sua maioria, requerem elevadas quantidades de energia durante sua operação, e são consideradas ambientalmente pouco sustentáveis. Em contraste, a produção biológica de hidrogênio ocorre sob condições próximas à temperatura e pressão ambientes, apresentando, portanto, menor demanda energética ao longo do processo \cite[authoryear][ref::das_0].

	Dentro das possibilidades de obtenção biológica de hidrogênio, encontram-se três rotas alternativas na natureza, cada uma com características que podem viabilizar ou inviabilizar sua aplicação, a depender das condições de operação e dos substratos utilizados: biofotólise da água, fotodecomposição e fermentação de compostos orgânicos \cite[authoryear][ref::das_0]. Essas são as principais alternativas para a biossíntese de hidrogênio, realizadas em ambientes anaeróbicos facultativos ou obrigatórios. Porém, tanto na biofotólise quanto na fotodecomposição, a formação de \chemical{H_{2}} é parcial ou completamente inibida na presença de oxigênio.

	A biofotólise adapta o processo de fotossíntese de plantas e algas para a geração de hidrogênio, em vez da produção de biomassa baseada em carbono \cite[authoryear][ref::das_0]. O processo ocorre por meio da absorção de luz, que promove a quebra da molécula de água em gás hidrogênio e gás oxigênio \cite[authoryear][ref::das_0]. Embora existam nuances relacionadas ao mecanismo completo da reação, de modo geral a formação de hidrogênio ocorre pela ação das enzimas hidrogenase ou nitrogenase, como ilustrado na \in{Figura}[fig:biophotolysis].

	\startplacefigure[title={Ilustração do mecanismo de biofotólise da água}, reference={fig:biophotolysis}]
		\externalfigure[image/water_photolysis.png][height=200pt]
	\stopplacefigure

	Na \in{Figura}[fig:biophotolysis], observa-se que a enzima (\chemical{Hidrogenase} ou \chemical{Nitrogenase}) desempenha um papel central na redução de prótons provenientes da fotólise da água. Os elétrons liberados nesse processo são transferidos até a enzima por meio da proteína \chemical{Ferredoxina}, na sua forma reduzida \chemical{Fd_{red}}. A oxidação da \chemical{Ferredoxina} pela \chemical{Hidrogenase} constitui a etapa final da cadeia de transferência eletrônica, na qual os \chemical{H^{+}} são reduzidos, resultando na formação de gás hidrogênio (\chemical{H_{2}}). Em seguida, a \chemical{Ferredoxina} é regenerada na forma oxidada (\chemical{Fd_{ox}}), sendo posteriormente novamente reduzida a \chemical{Fd_{red}}, o que permite a continuidade do ciclo redox. Observa-se também a formação de oxigênio molecular na etapa inicial da decomposição da água, conforme ilustrado na \in{Figura}[fig:biophotolysis]. Nesse contexto, o processo apresenta dependência de sistemas enzimáticos para viabilizar a redução de prótons a hidrogênio, caracterizando-o como um sistema indireto de produção de \chemical{H_{2}}. Em contrapartida, pode ocorrer a evolução direta de hidrogênio a partir da fotólise da água, configurando um sistema denominado direto \cite[authoryear][ref::sarangi_0]. Ambos os mecanismos podem ocorrer nesse tipo de rota metabólica; enquanto, o sistema direto apresenta maior simplicidade, o sistema indireto depende da atuação de enzimas específicas.

	A fotodecomposição de compostos orgânicos, também denominada fotofermentação, ocorre na presença de bactérias fotossintéticas \cite[authoryear][ref::sarangi_0], como as do genero \getbuffer[bac]. Esse processo caracteriza-se pela degradação de espécies químicas por ação da luz, com produção de hidrogênio molecular, de maneira análoga à biofotólise. A principal diferença entre os processos reside na natureza da fonte de alimentação e, consequentemente, na rota metabólica global envolvida, embora o princípio de conversão fotobiológica seja comparável. No caso da fotodecomposição, a matéria orgânica atua como substrato energético e doador de elétrons, permitindo maior flexibilidade espectral de absorção luminosa \cite[authoryear][ref::das_0]. Associado ao consumo de matéria orgânica, esse processo apresenta a vantagem de possibilitar a utilização de resíduos orgânicos como fonte de entrada, principalmente carboidratos \cite[authoryear][ref::ntaikou_0].

	\startplacefigure[title={Ilustração do mecanismo de fotodecomposiçao}, reference={fig:photo_fermentation}]
		\externalfigure[image/photo_fermentation.png][height=200pt]
	\stopplacefigure

	A \in{Figura}[fig:photo_fermentation] apresenta semelhanças com a \in{Figura}[fig:biophotolysis], destacando-se, em particular, a presença da \chemical{Ferredoxina}, a qual desempenha função equivalente em ambos os mecanismos. Esses dois sistemas diferem principalmente quanto à fonte de substrato, sendo que, no segundo caso, ácidos orgânicos são preferencialmente utilizados em substituição à água como doadores de elétrons. Outra analogia relevante pode ser estabelecida com o sistema da \in{Figura}[fig:alcoholic_fermentation], no qual o \chemical{NADH} desempenha papel funcional análogo ao da \chemical{Ferredoxina} nos processos de fotodecomposição e biofotólise da água. Em ambos os casos, esses cofatores atuam como transportadores de elétrons, viabilizando reações de redução na estequiometria final dos processos \cite[authoryear][ref::sarangi_0]. A \in{Equaçao}[eq:photo_fermentation] apresenta uma representação estequiométrica genérica do processo de fotodecomposição.

	\startplaceformula[reference={eq:photo_fermentation}]
		\startformula
			\startchemicalformula
				\chemical{{2}CH_{3}COOH} 
				\chemical{PLUS} 
				\chemical{H_{2}O} 
				\chemical{GIVES} 
				\chemical{CO_{2}} 
				\chemical{PLUS} 
				\chemical{{4}H_{2}}
			\stopchemicalformula
		\stopformula
	\stopplaceformula

	Na \in{Equação}[eq:photo_fermentation], observa-se a conversão do \chemical{Ácido acético} em \chemical{H_{2}}. Contudo, essa representação não explicita a dependência do processo em relação à enzima \chemical{Nitrogenase}, bem como à necessidade de \chemical{ATP} para sua atividade catalítica \cite[authoryear][ref::muller_0].

	A lista extendida das variações de substrato aceitas por sistemas de fotodecomposição de compostos orgânicos inclui monóxido de carbono em conjunto com água, em um processo denominado translação água-gás \cite[authoryear][ref::das_0], onde a estequiometria correspondente é apresentada na \in{Equação}[eq:shift].

	\startplaceformula[reference={eq:shift}]
		\startformula
			\startchemicalformula
				\chemical{CO} 
				\chemical{PLUS} 
				\chemical{H_{2}O} 
				\chemical{GIVES} 
				\chemical{CO_{2}} 
				\chemical{PLUS} 
				\chemical{H_{2}}
			\stopchemicalformula
		\stopformula
	\stopplaceformula

	A \in{Equação}[eq:shift] solidifica o conceito nomeado translação água-gás, ou {\emph water-gas shift}, que se encaixa no contexto da movimentação do oxigênio da água para o monóxido de carbono, oxidando o carbono, reduzindo o hidrogênio, e formando gás hidrogênio.

	Para ambos os sistemas envolvendo fotodecomposição, não há a formação de oxigênio molecular no meio reacional. Dessa forma, o planejamento do sistema é simplificado, uma vez que não há necessidade de controle rigoroso da presença de \chemical{O_2}, o qual pode atuar como um forte inibidor da atividade de hidrogenases. O efeito inibitório do oxigênio está associado à sensibilidade dessas enzimas ao meio oxidante, podendo ocorrer inativação reversível ou irreversível do sítio ativo contendo metais de transição, particularmente centros \chemical{Fe-Fe} ou \chemical{Ni-Fe}. Nessas condições, a presença de \chemical{O_{2}} leva à oxidação desses centros catalíticos, reduzindo significativamente a eficiência da redução de prótons a \chemical{H_{2}} (CITAÇAO). Em sistemas como a biofotólise da água, essa limitação exige estratégias adicionais de separação temporal ou espacial entre a etapa de geração de oxigênio e a etapa de produção de hidrogênio, a fim de preservar a atividade enzimática.

	A fermentação de hidrogênio constitui o último dos três principais métodos biológicos de produção de \chemical{H_{2}}, apresentando cinética e condições reacionais distintas, o que resulta em características sistemáticas particulares. Bactérias fermentativas podem apresentar elevadas taxas de evolução de hidrogênio, com capacidade de produção tanto na presença quanto na ausência de luz, não havendo dependência de iluminação, o que permite sua atividade contínua em diferentes ciclos ambientais. Além da produção de \chemical{H_{2}} em quantidades relativamente elevadas quando comparadas a outros sistemas biológicos, esses microrganismos geralmente apresentam taxas de crescimento adequadas para aplicações em processos produtivos, o que pode favorecer o aumento da conversão sob condições ambientais otimas \cite[authoryear][ref::das_0]. Também denominada fermentação escura, esse processo pode ocorrer em condições estritamente anaeróbias ou facultativamente anaeróbias, dependendo do microrganismo envolvido, com uma estequiometria genérica apresentada na \in{Equação}[eq:dark_fermentation]. Exemplos incluem bactérias do gênero \getbuffer[clos] \cite[authoryear][ref::muller_0].

	\startplaceformula[reference={eq:dark_fermentation}]
		\startformula
			\startchemicalformula
				\chemical{C_{6}H_{12}O_{6}}
				\chemical{PLUS}
				\chemical{{2}H_{2}O}
				\chemical{GIVES}
				\chemical{{2}CH_{3}COOH}
				\chemical{PLUS}
				\chemical{CO_{2}}
				\chemical{PLUS}
				\chemical{{4}H_{2}}
			\stopchemicalformula
		\stopformula
	\stopplaceformula

	A \in{Equação}[eq:dark_fermentation] exemplifica uma das possíveis reações globais realizadas por microrganismos fermentativos anaeróbios heterotróficos, na qual ocorre o consumo de \chemical{Glicose} com formação de \chemical{ácido acético} e \chemical{H_{2}}. O \chemical{ácido acético} é um ácido orgânico que pode ser posteriormente consumido por microrganismos fotodecompositores, de modo que a produção global de hidrogênio é potencializada pela combinação entre fermentação escura e fotofermentação \cite[authoryear][ref::muller_0]. Dessa forma, um arranjo em série das reações representadas nas \in{Equação}[eq:dark_fermentation] e \in{Equação}[eq:photo_fermentation] permite maior aproveitamento dos ácidos orgânicos intermediários, resultando em conversão globais mais eficientes de substrato em \chemical{H_{2}}. Outro arranjo possivel é o reaproveitamento dos subprodutos da fermentação escura na eletrólise microbiana. Esta possui uma estequimetria semelhante a fotofermentação, com distinções em torno da força motriz, mudando para oxiredução em eletrodos.
\stopsection

